% ===========================================================================
% CHƯƠNG 2: XÂY DỰNG MÔ HÌNH HỌC MÁY TRONG HỖ TRỢ KIỂM DUYỆT TÍN DỤNG
% ===========================================================================
\chapter{XÂY DỰNG MÔ HÌNH HỌC MÁY TRONG HỖ TRỢ KIỂM DUYỆT TÍN DỤNG}
% Tiêu đề chạy rút gọn: tiêu đề chương đầy đủ dài hơn một dòng đầu trang,
% khiến fancyhdr phải nới \headheight và làm lệch khung chữ của cả chương.
\markboth{Chương 2. Xây dựng mô hình học máy hỗ trợ kiểm duyệt tín dụng}{}
\label{chap:xay-dung-he-thong}

\noindent\textit{Chương 2 trình bày thực trạng quy trình thẩm định tín dụng, xác
định bài toán học máy và phương pháp thu thập, tiền xử lý dữ liệu hồ sơ vay. Nội
dung chương tập trung vào kỹ thuật trích xuất đặc trưng, gán nhãn rủi ro và thiết
lập quy trình huấn luyện, tối ưu mô hình làm tiền đề cho các thực nghiệm trình
bày ở Chương~\ref{chap:thuc-nghiem}.}

\section{Phân tích bài toán kiểm duyệt tín dụng}
\label{sec:phan-tich-bai-toan}

\subsection{Xác định các mức độ rủi ro và trạng thái khoản vay cần dự báo}
\label{subsec:muc-do-rui-ro}

\paragraph{Biến mục tiêu.} Đối tượng dự báo của mô hình là xác suất một hồ sơ
vay phát sinh nợ xấu trong cửa sổ quan sát kể từ thời điểm cấp tín dụng. Biến mục
tiêu được định nghĩa nhị phân:

\begin{equation}
y_i =
\begin{cases}
1, & \text{nếu hồ sơ } i \text{ phát sinh nợ nhóm 3--5 (quá hạn} \ge 90 \text{ ngày)} \\
   & \text{trong vòng 12 tháng kể từ ngày cấp tín dụng,} \\[4pt]
0, & \text{trong trường hợp còn lại.}
\end{cases}
\label{eq:target}
\end{equation}

\paragraph{Ba trạng thái thay vì hai.} Mặc dù biến mục tiêu là nhị phân, quá
trình gán nhãn phải phân biệt \textbf{ba} trạng thái. Ngoài hai trạng thái
\textit{Tốt} và \textit{Xấu}, cần nhận diện trạng thái \textit{Chưa xác định}
dành cho các hồ sơ chưa đi hết cửa sổ quan sát tại thời điểm trích xuất dữ liệu.
Nếu bỏ qua và mặc định xếp nhóm này vào loại Tốt, dữ liệu huấn luyện sẽ chứa nhãn
sai một cách có hệ thống ở chính những hồ sơ mới nhất -- nhóm gần nhất với điều
kiện hiện tại. Hệ quả là mô hình đánh giá thấp rủi ro của các hồ sơ gần đây. Đề
án loại nhóm Chưa xác định khỏi tập huấn luyện và tập đánh giá.

\paragraph{Thang phân hạng rủi ro.} Kết quả đầu ra không dừng ở một giá trị xác
suất. Để phục vụ nghiệp vụ, thang $\PD$ liên tục được ánh xạ sang bảy hạng rủi ro
theo thông lệ xếp hạng tín dụng, trình bày tại Bảng~\ref{tab:risk-bands-design}.

\begin{table}[htbp]
\centering
\caption{Thang phân hạng rủi ro theo xác suất vỡ nợ}
\label{tab:risk-bands-design}
\small
\begin{tabular}{@{}cccl@{}}
\toprule
\textbf{Hạng} & \textbf{Khoảng $\PD$} & \textbf{Mức rủi ro} &
\textbf{Định hướng xử lý} \\
\midrule
AA  & $[0{,}00;\ 0{,}02)$ & Rất thấp    & Ưu tiên duyệt, có thể ưu đãi lãi suất \\
A   & $[0{,}02;\ 0{,}05)$ & Thấp        & Duyệt theo điều kiện chuẩn \\
BBB & $[0{,}05;\ 0{,}08)$ & Trung bình thấp & Duyệt, có thể yêu cầu bổ sung điều kiện \\
BB  & $[0{,}08;\ 0{,}15)$ & Trung bình  & Chuyển thẩm định thủ công \\
B   & $[0{,}15;\ 0{,}25)$ & Trung bình cao & Chuyển thẩm định thủ công \\
CCC & $[0{,}25;\ 0{,}40)$ & Cao         & Thẩm định kỹ hoặc từ chối \\
D   & $[0{,}40;\ 1{,}00]$ & Rất cao     & Từ chối \\
\bottomrule
\end{tabular}
\end{table}

\subsection{Đặc trưng đầu vào và đầu ra của mô hình kiểm duyệt tín dụng}
\label{subsec:dau-vao-dau-ra}

\paragraph{Đầu vào.} Vectơ đặc trưng của mỗi hồ sơ được tổng hợp từ bốn nhóm
nguồn dữ liệu đã phân tích ở mục~\ref{subsec:nhom-du-lieu}, bổ sung thêm nhóm thứ
năm gồm các biến tổ hợp và biến tương tác. Nguyên tắc bắt buộc đối với mọi đặc
trưng là \textbf{tính đúng đắn theo thời điểm}: giá trị của mọi biến chỉ được
tính từ dữ liệu tồn tại trước hoặc tại thời điểm nộp hồ sơ.

\paragraph{Đầu ra.} Khác với bài toán phân loại thông thường chỉ trả về một nhãn,
kết quả đầu ra của mô hình gồm năm thành phần, phản ánh vai trò hỗ trợ chứ không
thay thế con người:

\begin{enumerate}[label=(\arabic*)]
  \item \textbf{Xác suất vỡ nợ} $\PD \in (0, 1)$ đã được hiệu chuẩn;
  \item \textbf{Hạng rủi ro} theo thang bảy mức ở Bảng~\ref{tab:risk-bands-design};
  \item \textbf{Khuyến nghị} thuộc một trong ba giá trị: \code{APPROVE} (khuyến
  nghị duyệt), \code{REVIEW} (chuyển thẩm định thủ công), \code{REJECT} (khuyến
  nghị từ chối);
  \item \textbf{Danh sách lý do} (\eng{reason codes}) gồm các luật nghiệp vụ bị
  vi phạm và các chỉ tiêu định lượng có ảnh hưởng lớn nhất tới kết quả;
  \item \textbf{Đề xuất điều kiện cấp tín dụng} gồm hạn mức, kỳ hạn và mức lãi
  suất theo rủi ro, trong trường hợp khuyến nghị duyệt.
\end{enumerate}

Thành phần thứ tư (danh sách lý do) là điểm phân biệt giữa mô hình đóng vai trò
hỗ trợ và mô hình ra quyết định tự động. Thiếu thành phần này, cán bộ tín dụng
chỉ nhận được
một con số mà không có căn cứ để chấp nhận hay bác bỏ, dẫn tới hai thái cực đều
không mong muốn: chấp nhận kết quả một cách thụ động hoặc bỏ qua hoàn toàn khuyến
nghị của mô hình.

\subsection{Yêu cầu kỹ thuật và pháp lý đối với mô hình kiểm duyệt tín dụng}
\label{subsec:yeu-cau}

\subsubsection{Yêu cầu chức năng}

Quy trình kiểm duyệt phải thực hiện được tám chức năng, tương ứng các chặng của
một lượt xử lý hồ sơ. Trên \textit{đường xử lý chính}, quy trình
\textbf{tiếp nhận hồ sơ} qua
giao diện dịch vụ và kiểm tra tính đầy đủ, hợp lệ của dữ liệu đầu vào;
\textbf{tổng hợp dữ liệu} khách hàng, tài khoản, giao dịch và bản tin CIC về cấp
độ hồ sơ; \textbf{xây dựng đặc trưng} theo đúng định nghĩa đã dùng khi huấn
luyện, bảo đảm tính đúng đắn theo thời điểm; \textbf{chấm điểm rủi ro} bằng mô
hình đã hiệu chuẩn; \textbf{áp dụng bộ luật nghiệp vụ}; và \textbf{đưa ra khuyến
nghị} ba mức theo thứ tự ưu tiên xác định, kèm \textbf{danh sách lý do} và
\textbf{đề xuất điều kiện cấp tín dụng} (hạn mức, kỳ hạn, lãi suất theo rủi ro).

Ngoài đường xử lý chính, quy trình cần hỗ trợ \textbf{chấm điểm theo lô} phục vụ
tái đánh giá danh mục hiện hữu, và \textbf{giám sát mô hình} thông qua theo dõi
trôi dạt dữ liệu cùng chất lượng dự báo theo thời gian.

\subsubsection{Yêu cầu phi chức năng}

\begin{itemize}
  \item \textbf{Tính nhất quán:} cùng một hồ sơ với cùng dữ liệu đầu vào phải
  luôn cho cùng kết quả. Yêu cầu này đòi hỏi định nghĩa đặc trưng khi huấn luyện
  và khi chấm điểm phải trùng khớp tuyệt đối.

  \item \textbf{Thời gian phản hồi:} thời gian chấm điểm một hồ sơ phải đủ nhỏ để
  tích hợp được vào giao diện làm việc trực tuyến của cán bộ tín dụng. Đề án đặt
  mục tiêu thời gian xử lý một hồ sơ ở mức dưới 100 ms.

  \item \textbf{Khả năng hoạt động khi thiếu dữ liệu:} mô hình phải trả về kết
  quả cả khi không truy vấn được bản tin CIC. Trong trường hợp đó, kết quả là
  \code{REVIEW} kèm lý do, chứ không phải một thông báo lỗi.

  \item \textbf{Khả năng tái lập:} toàn bộ quy trình từ dữ liệu tới kết quả phải
  tái lập được với cùng một hạt giống ngẫu nhiên.
\end{itemize}

\subsubsection{Yêu cầu pháp lý và tuân thủ}

Mô hình sử dụng trong hoạt động cấp tín dụng chịu ràng buộc bởi các quy định pháp
lý, trong đó năm nhóm yêu cầu sau ảnh hưởng trực tiếp tới thiết kế.

\textit{Giải trình được:} mọi khuyến nghị phải kèm lý do có thể trình bày cho
khách hàng và cơ quan quản lý. Yêu cầu này loại trừ khả năng sử dụng mô hình như
một hộp đen, và là lý do đề án tích hợp cơ chế giải thích SHAP ngay trong thiết
kế chứ không coi đó là phần bổ sung.

\textit{Không phân biệt đối xử:} mô hình không được sử dụng các biến bị cấm hoặc
biến đại diện cho chúng làm căn cứ phân biệt. Trong tập đặc trưng của đề án, biến
giới tính được giữ lại phục vụ phân tích thống kê mô tả nhưng cần cân nhắc loại
bỏ khi triển khai thực tế, tùy quy định áp dụng.

\textit{Tuân thủ định nghĩa phân loại nợ:} nhãn rủi ro phải được xây dựng theo
đúng quy định hiện hành~\cite{tt11}, bảo đảm kết quả của mô hình tương thích với
các báo cáo quản trị rủi ro và trích lập dự phòng.

\textit{Bảo mật dữ liệu cá nhân:} đề án sử dụng dữ liệu mô phỏng; khi áp dụng
trên dữ liệu thực, dữ liệu cần được ẩn danh và giới hạn quyền truy
cập~\cite{tt15}.

\textit{Thẩm quyền quyết định thuộc về con người:} đây là ràng buộc bao trùm --
mô hình cung cấp thông tin hỗ trợ, quyết định cấp tín dụng cuối cùng thuộc về
người hoặc bộ phận có thẩm quyền theo phân cấp của tổ chức tín dụng.

\section{Thu thập và tổng hợp dữ liệu}
\label{sec:thu-thap-du-lieu}

\subsection{Nguồn dữ liệu tín dụng truyền thống và dữ liệu bổ sung}
\label{subsec:nguon-du-lieu}

Hệ thống sử dụng năm bảng dữ liệu, chia thành hai nhóm.

\paragraph{Dữ liệu tín dụng truyền thống.} Gồm ba nguồn đã được sử dụng từ lâu
trong thẩm định tín dụng: (i) \textit{hồ sơ khách hàng} chứa thông tin nhân khẩu
học, nghề nghiệp và thu nhập khai báo; (ii) \textit{đơn đề nghị vay} chứa tham số
khoản vay; và (iii) \textit{bản tin CIC} chứa lịch sử quan hệ tín dụng của khách
hàng với toàn hệ thống tổ chức tín dụng.

\paragraph{Dữ liệu bổ sung.} Gồm hai nguồn mà tổ chức tín dụng tự sở hữu nhưng
thường chưa được khai thác đầy đủ trong thẩm định: (iv) \textit{thông tin tài
khoản} và (v) \textit{lịch sử giao dịch}. Nhóm này có ba ưu thế: không phụ thuộc
dịch vụ bên ngoài, phản ánh tình trạng tài chính hiện tại thay vì lịch sử, và đặc
biệt có giá trị với nhóm khách hàng thiếu lịch sử tín dụng -- nhóm mà bản tin CIC
gần như không cung cấp thông tin.

Bảng~\ref{tab:data-schema} mô tả lược đồ của năm nguồn dữ liệu.

\begin{table}[htbp]
\centering
\caption{Lược đồ các bảng dữ liệu sử dụng trong đề án}
\label{tab:data-schema}
\small
\begin{tabular}{@{}L{2.4cm}L{2.7cm}L{7.0cm}@{}}
\toprule
\textbf{Bảng} & \textbf{Khóa chính} & \textbf{Nội dung} \\
\midrule
\code{customers} & \code{customer\_id} &
Nhân khẩu học, nghề nghiệp, thu nhập khai báo, nơi cư trú và thời gian quan hệ
với ngân hàng \\[2pt]
\code{accounts} & \code{account\_id} &
Loại tài khoản, ngày mở, số dư bình quân ba tháng, số dư hiện tại, trạng thái
\\[2pt]
\code{transactions} & \code{txn\_id} &
Thời điểm, chiều giao dịch, số tiền, nhóm giao dịch và kênh thực hiện \\[2pt]
\code{applications} & \code{application\_id} &
Thời điểm nộp, sản phẩm, số tiền đề nghị, kỳ hạn, lãi suất, mục đích vay, kênh
tiếp nhận, tài sản bảo đảm và khoản trả góp \\[2pt]
\code{cic\_reports} & \code{report\_id} &
Thời điểm truy vấn, điểm CIC, nhóm nợ, số ngày quá hạn, số tổ chức có quan hệ, dư
nợ, hạn mức và số lần truy vấn tín dụng \\
\bottomrule
\end{tabular}
\end{table}

\subsection{Thu thập và trích xuất dữ liệu khách hàng vay vốn}
\label{subsec:trich-xuat-du-lieu}

\paragraph{Đơn vị phân tích.} Toàn bộ dữ liệu được quy về đơn vị phân tích là
\textit{hồ sơ vay}: mỗi dòng trong ma trận đặc trưng cuối cùng tương ứng một hồ
sơ, được xác định bởi cặp (mã hồ sơ, thời điểm nộp). Việc chọn hồ sơ -- thay vì
khách hàng -- làm đơn vị phân tích là cần thiết vì một khách hàng có thể nộp
nhiều hồ sơ ở các thời điểm khác nhau, và tình trạng tài chính tại mỗi thời điểm
là khác nhau.

\paragraph{Quy tắc ghép nối theo thời điểm.} Đây là điểm kỹ thuật quan trọng nhất
của bước thu thập dữ liệu. Với mỗi hồ sơ nộp tại thời điểm $t$, việc ghép nối các
nguồn phải tuân thủ:

\begin{itemize}
  \item Thông tin khách hàng và tài khoản: lấy trạng thái tại thời điểm $t$;
  \item Bản tin CIC: chỉ sử dụng bản tin được truy vấn tại thời điểm không muộn
  hơn $t$;
  \item Giao dịch: chỉ tổng hợp các giao dịch có thời điểm phát sinh nhỏ hơn hoặc
  bằng $t$.
\end{itemize}

Vi phạm bất kỳ quy tắc nào trong ba quy tắc trên đều dẫn tới rò rỉ thông tin từ
tương lai. Đề án đưa quy tắc thứ nhất và thứ ba vào chính thuật toán tính đặc
trưng (mục~\ref{subsec:dac-trung-hanh-vi}), còn quy tắc thứ hai được kiểm tra
tường minh trong bước kiểm tra tính hợp lệ dữ liệu.

\subsection{Phân tích thống kê mô tả tập dữ liệu tín dụng}
\label{subsec:thong-ke-mo-ta}

Bước thống kê mô tả lập bảng đặc tả cho từng biến, gồm: kiểu dữ liệu, số quan sát
hợp lệ, tỷ lệ giá trị khuyết, số giá trị phân biệt, và với biến định lượng là các
phân vị 1\%, 25\%, 50\%, 75\%, 99\% cùng trung bình và độ lệch chuẩn.

Bảng mô tả phục vụ ba quyết định ở các bước sau. Tỷ lệ khuyết của từng biến quyết
định chiến lược điền khuyết và xác định biến nào cần loại bỏ do độ phủ quá thấp.
Khoảng cách giữa phân vị 99\% và giá trị lớn nhất cho biết mức độ nặng của đuôi
phân phối, làm căn cứ quyết định có cần cắt ngọn hay không. Số giá trị phân biệt
của biến định tính quyết định cách mã hóa: biến có ít hạng mục phù hợp với mã hóa
one-hot, biến có nhiều hạng mục cần cơ chế xử lý riêng.

\subsection{Phát hiện và xử lý dữ liệu ngoại lệ, dữ liệu không hợp lệ}
\label{subsec:du-lieu-khong-hop-le}

Hai khái niệm dưới đây thường bị gộp chung nhưng có bản chất khác nhau và đòi
hỏi cách xử lý khác nhau.

\paragraph{Dữ liệu không hợp lệ} là dữ liệu \textit{sai} -- vi phạm ràng buộc
nghiệp vụ hoặc ràng buộc logic. Ví dụ: tuổi khách hàng bằng 15, thu nhập âm, nhóm
nợ xấu nhất trong 12 tháng lại tốt hơn nhóm nợ hiện tại, hoặc bản tin CIC được
truy vấn sau thời điểm nộp hồ sơ. Đây là lỗi dữ liệu và phải được loại bỏ hoặc
sửa chữa.

\paragraph{Giá trị ngoại lệ} là giá trị nằm xa phần lớn quan sát nhưng \textit{có
thể đúng}. Một khách hàng có thu nhập gấp hai mươi lần trung vị hoặc một giao
dịch trị giá vài tỷ đồng là bất thường về mặt thống kê nhưng hoàn toàn có thể có
thật. Loại bỏ những quan sát này đồng nghĩa với việc loại bỏ chính nhóm khách
hàng mà tổ chức tín dụng quan tâm.

Đề án xây dựng bộ mười lăm ràng buộc nghiệp vụ để phát hiện dữ liệu không hợp lệ,
chia theo ba bảng dữ liệu chính. Bảng~\ref{tab:validity-rules} trình bày các ràng
buộc này.

\begin{table}[htbp]
\centering
\caption{Bộ ràng buộc kiểm tra tính hợp lệ của dữ liệu}
\label{tab:validity-rules}
\small
\setlength{\tabcolsep}{5pt}
\renewcommand{\arraystretch}{1.15}
\begin{tabular}{@{}L{2.5cm}L{4.8cm}L{6.6cm}@{}}
\toprule
\textbf{Bảng} & \textbf{Mã ràng buộc} & \textbf{Nội dung} \\
\midrule
\code{customers}
  & \code{CUS\_AGE\_RANGE}        & Tuổi khách hàng trong khoảng 18--75 \\
  & \code{CUS\_INCOME\_POSITIVE}  & Thu nhập khai báo lớn hơn 0 \\
  & \code{CUS\_TENURE\_NONNEG}    & Thâm niên công việc không âm \\
  & \code{CUS\_DEPENDENTS\_RANGE} & Số người phụ thuộc trong khoảng 0--10 \\
\midrule
\code{applications}
  & \code{APP\_AMOUNT\_POSITIVE}     & Số tiền đề nghị vay lớn hơn 0 \\
  & \code{APP\_TENOR\_RANGE}         & Kỳ hạn vay trong khoảng 3--360 tháng \\
  & \code{APP\_RATE\_RANGE}          & Lãi suất trong khoảng 0--40\%/năm \\
  & \code{APP\_INSTALLMENT\_POSITIVE} & Khoản trả góp hàng tháng lớn hơn 0 \\
  & \code{APP\_LTV\_RANGE}           & Tỷ lệ vay trên giá trị tài sản bảo đảm
  không vượt 1,2 \\
\midrule
\code{cic\_reports}
  & \code{CIC\_SCORE\_RANGE}       & Điểm CIC trong khoảng 300--850 \\
  & \code{CIC\_GROUP\_RANGE}       & Nhóm nợ trong khoảng 1--5 \\
  & \code{CIC\_GROUP\_ORDER}       & Nhóm nợ xấu nhất 12 tháng không tốt hơn
  nhóm nợ hiện tại \\
  & \code{CIC\_DPD\_NONNEG}        & Số ngày quá hạn không âm \\
  & \code{CIC\_UTILIZATION\_RANGE} & Tỷ lệ sử dụng hạn mức trong khoảng 0--2 \\
  & \code{CIC\_QUERY\_BEFORE\_APPLY} & Bản tin CIC được truy vấn không muộn hơn
  thời điểm nộp hồ sơ \\
\bottomrule
\end{tabular}
\end{table}

\paragraph{Kiểm chứng cơ chế phát hiện.} Một bảng kết quả kiểm tra không ghi nhận
vi phạm nào có thể tương ứng hai khả năng: dữ liệu thực sự sạch, hoặc cơ chế kiểm
tra không hoạt động. Để phân biệt, đề án thực hiện phép \textbf{tiêm lỗi}: tạo
một bản sao dữ liệu, cố ý đưa vào một tỷ lệ lỗi đã biết trước ứng với từng ràng
buộc, rồi chạy lại cơ chế kiểm tra và đối chiếu số lỗi phát hiện được với số lỗi
đã tiêm. Kết quả kiểm chứng được trình bày ở
mục~\ref{subsec:kq-chat-luong-du-lieu}.

\paragraph{Xử lý giá trị ngoại lệ.} Đối với giá trị ngoại lệ, đề án không loại bỏ
quan sát mà áp dụng kỹ thuật \textbf{cắt ngọn} (\eng{winsorize}) tại phân vị 1\%
và 99\% học từ tập huấn luyện. Cách này giữ nguyên số quan sát -- do đó không làm
mất nhóm khách hàng đặc thù -- đồng thời hạn chế ảnh hưởng quá mức của vài giá
trị cực trị lên các mô hình nhạy cảm với thang đo như hồi quy logistic, SVM và
mạng nơ-ron. Việc phát hiện ngoại lệ sử dụng quy tắc $k \times \text{IQR}$ với
$k = 3$ thay vì giá trị thông thường $k = 1{,}5$, bởi các biến tiền tệ trong dữ
liệu tín dụng có phân phối lệch phải mạnh; ngưỡng $1{,}5 \times \text{IQR}$ sẽ
đánh dấu một tỷ lệ rất lớn quan sát là ngoại lệ mà không mang lại thông tin hữu
ích.

\subsection{Phân chia dữ liệu theo khoảng thời gian và nhóm khách hàng}
\label{subsec:phan-chia-du-lieu}

Đề án sử dụng hai chiến lược phân chia với hai mục đích khác nhau.

\paragraph{Phân chia theo thời gian (chính).} Dữ liệu được sắp xếp theo thời điểm
nộp hồ sơ và chia thành ba đoạn liên tiếp: 70\% đầu làm tập huấn luyện, 15\% tiếp
theo làm tập kiểm định, 15\% cuối làm tập kiểm thử. Chiến lược này mô phỏng đúng
điều kiện vận hành thực tế -- mô hình luôn được huấn luyện trên quá khứ và áp
dụng cho tương lai -- nên phản ánh cả ảnh hưởng của trôi dạt dữ liệu.

\paragraph{Phân chia theo nhóm khách hàng (phương án bổ sung).} Toàn bộ hồ sơ của
một khách hàng được đặt vào cùng một tập. Chiến lược này loại trừ rò rỉ thông tin
theo khách hàng nhưng không phản ánh trôi dạt theo thời gian, nên chỉ phù hợp khi
tập dữ liệu có nhiều khách hàng nộp lặp nhiều hồ sơ. Hệ thống hiện thực sẵn
phương án này để có thể đối chiếu khi áp dụng trên dữ liệu thực; toàn bộ kết quả
thực nghiệm trình bày ở Chương~\ref{chap:thuc-nghiem} sử dụng cách chia theo thời
gian.

\paragraph{Tách riêng vai trò trong tập huấn luyện.} Với các mô hình có cơ chế
dừng sớm (gradient boosting), cần một tập dữ liệu độc lập để đánh giá tiêu chí
dừng. Nếu lấy tập này từ tập kiểm định, tập kiểm định sẽ đồng thời phục vụ hai
mục đích -- dừng sớm và hiệu chuẩn xác suất -- làm suy giảm tính độc lập của cả
hai. Đề án tách 15\% cuối của tập huấn luyện thành tập dừng sớm riêng biệt, cho
cấu trúc bốn phần:

\begin{center}
\begin{tabular}{@{}cccc@{}}
\toprule
\textbf{Tập} & \textbf{Tỷ trọng} & \textbf{Vai trò} & \textbf{Sử dụng} \\
\midrule
\code{train\_fit} & 59,5\% & Huấn luyện & Mọi mô hình \\
\code{train\_es}  & 10,5\% & Dừng sớm   & Mô hình boosting \\
\code{valid}      & 15,0\% & Hiệu chuẩn xác suất, chọn ngưỡng & Mọi mô hình \\
\code{test}       & 15,0\% & Đánh giá cuối cùng & Chỉ dùng một lần \\
\bottomrule
\end{tabular}
\end{center}

Nguyên tắc quan trọng nhất là tập \code{test} chỉ được sử dụng \textit{một lần
duy nhất} ở bước báo cáo kết quả cuối cùng. Mọi quyết định về lựa chọn mô hình,
siêu tham số và ngưỡng đều dựa trên tập \code{valid}.

\section{Tiền xử lý và làm sạch dữ liệu}
\label{sec:tien-xu-ly}

Nguyên tắc xuyên suốt của toàn bộ khối tiền xử lý: \textbf{mọi tham số biến đổi
chỉ được ước lượng trên tập huấn luyện} rồi áp dụng cho các tập còn lại. Nếu ước
lượng trung vị, trung bình hay phân vị trên toàn bộ dữ liệu, thông tin từ tập
kiểm thử rò rỉ vào quá trình huấn luyện và kết quả đánh giá sẽ lạc quan hơn thực
tế.

\subsection{Xử lý giá trị khuyết thiếu}
\label{subsec:xu-ly-khuyet}

Trong dữ liệu tín dụng, giá trị khuyết không đồng nghĩa với lỗi thu thập. Đề án
phân biệt hai loại và xử lý theo hai cách khác nhau.

\paragraph{Khuyết mang thông tin.} Khách hàng chưa từng có quan hệ tín dụng
(\eng{thin-file}) sẽ có toàn bộ trường CIC trống. Bản thân trạng thái ``không có
lịch sử tín dụng'' là một tín hiệu rủi ro riêng biệt. Cách xử lý đúng gồm hai
phần: điền giá trị theo \textit{quy ước nghiệp vụ} và bổ sung \textit{biến chỉ
báo khuyết} để mô hình tự học mức rủi ro thực tế của nhóm này. Quy ước điền được
chọn theo hướng trung tính -- không giả định khách hàng tốt cũng không giả định
khách hàng xấu: điểm CIC đặt bằng 600 (mức giữa của thang 300--850), nhóm nợ đặt
bằng 1 (chưa phát sinh nợ), các chỉ tiêu đếm và dư nợ đặt bằng 0.

\paragraph{Khuyết ngẫu nhiên.} Với các trường như thu nhập khai báo hay thâm niên
công việc, giá trị khuyết phản ánh việc khách hàng không cung cấp thông tin. Đề
án điền bằng \textit{trung vị của tập huấn luyện} kèm biến chỉ báo. Trung vị được
chọn thay vì trung bình vì ít nhạy cảm với đuôi phân phối lệch phải điển hình của
các biến tiền tệ.

Việc điền theo thống kê được thực hiện \textit{bên trong} pipeline huấn luyện của
scikit-learn, nhờ đó tham số điền được ước lượng lại ở mỗi lần khớp mô hình trong
kiểm định chéo, loại bỏ hoàn toàn nguy cơ rò rỉ.

\subsection{Phát hiện và xử lý dữ liệu ngoại lệ}
\label{subsec:xu-ly-ngoai-le}

Quy trình xử lý ngoại lệ gồm hai bước. Bước thứ nhất -- thống kê -- sử dụng quy
tắc $3 \times \text{IQR}$ để lập bảng các biến có giá trị nằm ngoài ngưỡng, kèm
tỷ lệ và giá trị các phân vị biên. Bảng này là cơ sở để đánh giá mức độ lệch của
từng biến.

Bước thứ hai -- xử lý -- áp dụng phép cắt ngọn tại phân vị 1\% và 99\%. Ngưỡng
cắt được học từ tập huấn luyện và áp dụng nguyên vẹn cho tập kiểm định và tập
kiểm thử; nếu học lại trên từng tập, phép biến đổi sẽ khác nhau giữa các tập và
làm sai lệch kết quả đánh giá.

Phép cắt ngọn chỉ cần thiết với nhóm mô hình nhạy cảm với thang đo. Các mô hình dựa trên cây chia theo ngưỡng nên bất biến với phép biến
đổi đơn điệu, do đó không chịu ảnh hưởng của giá trị cực trị theo cách tương tự.

\subsection{Chuẩn hóa và mã hóa biến định tính/định lượng}
\label{subsec:chuan-hoa-ma-hoa}

Mỗi họ mô hình có yêu cầu khác nhau về dạng đầu vào, nên đề án sử dụng bốn phương
án biến đổi song song thay vì một phương án chung
(Bảng~\ref{tab:encoding-schemes}).

\begin{table}[htbp]
\centering
\caption{Bốn phương án biến đổi dữ liệu theo yêu cầu của từng họ mô hình}
\label{tab:encoding-schemes}
\small
\begin{tabular}{@{}L{2.3cm}L{4.4cm}L{5.5cm}@{}}
\toprule
\textbf{Phương án} & \textbf{Nội dung biến đổi} & \textbf{Mô hình áp dụng và lý
do} \\
\midrule
\code{numeric\_ohe} &
Điền khuyết bằng trung vị, chuẩn hóa z-score cho biến định lượng; mã hóa one-hot
cho biến định tính &
Hồi quy logistic, SVM, mạng nơ-ron. Các mô hình dựa trên khoảng cách hoặc
gradient nhạy cảm với chênh lệch thang đo giữa các biến \\[2pt]
\code{ordinal} &
Điền khuyết; mã hóa thứ tự cho biến định tính &
Cây quyết định, rừng ngẫu nhiên, Extra Trees, Hist Gradient Boosting. Cây chia
theo ngưỡng nên bất biến với phép biến đổi đơn điệu; chuẩn hóa không cần thiết và
làm mất khả năng diễn giải ngưỡng \\[2pt]
\code{native\_cat} &
Giữ nguyên biến định tính ở dạng hạng mục &
XGBoost, LightGBM, CatBoost. Các thư viện này có cơ chế xử lý biến định tính
riêng, hiệu quả hơn mã hóa bên ngoài \\[2pt]
\code{woe} &
Rời rạc hóa và thay mỗi khoảng bằng giá trị $\WOE$ &
Mô hình thẻ điểm. Đưa mọi biến về thang log-odds và cho phép áp ràng buộc đơn
điệu \\
\bottomrule
\end{tabular}
\end{table}

Toàn bộ phép biến đổi được gói trong cấu trúc \code{ColumnTransformer} của
scikit-learn, bảo đảm tham số chỉ được ước lượng trên dữ liệu huấn luyện của từng
lần khớp mô hình.

\paragraph{Xử lý hạng mục chưa từng xuất hiện.} Khi chấm điểm một hồ sơ mới, biến
định tính có thể nhận giá trị chưa từng xuất hiện trong tập huấn luyện. Đề án
``đóng băng'' danh mục hạng mục theo tập huấn luyện và lưu kèm mô hình; hạng mục
lạ được quy về giá trị khuyết. Đây là hành vi đúng khi triển khai: mô hình không
có cơ sở nào để đánh giá một hạng mục chưa từng quan sát.

\subsection{Phân chia tập dữ liệu huấn luyện, kiểm thử}
\label{subsec:chia-tap}

Cấu trúc bốn tập dữ liệu đã được trình bày ở mục~\ref{subsec:phan-chia-du-lieu}.
Nội dung cần làm rõ ở đây là \textbf{thứ tự thao tác} giữa việc phân chia và các
phép tiền xử lý:

\begin{center}
gán nhãn $\rightarrow$ loại hồ sơ chưa đủ cửa sổ quan sát $\rightarrow$ chia tập
$\rightarrow$ học tham số biến đổi trên tập huấn luyện $\rightarrow$ áp dụng cho
các tập còn lại.
\end{center}

Việc phân chia phải diễn ra \textit{trước} mọi phép biến đổi có ước lượng tham
số. Nếu đảo ngược, chẳng hạn chuẩn hóa trên toàn bộ dữ liệu rồi mới chia tập,
trung bình và độ lệch chuẩn dùng để chuẩn hóa sẽ chứa thông tin của tập kiểm thử;
mô hình khi đó được đánh giá trên dữ liệu mà nó đã gián tiếp tiếp xúc, cho kết
quả lạc quan hơn thực tế.

\section{Trích xuất và lựa chọn đặc trưng}
\label{sec:trich-xuat-dac-trung}

\subsection{Trích xuất nhóm đặc trưng thông tin cá nhân và tài chính}
\label{subsec:dac-trung-ca-nhan}

Nhóm này gồm mười lăm đặc trưng lấy trực tiếp từ hồ sơ khách hàng và thông tin
tài khoản, chia thành ba mảng: \textit{nhân khẩu học} (tuổi, giới tính, tình
trạng hôn nhân, số người phụ thuộc, trình độ học vấn, tỉnh/thành cư trú, hình
thức nhà ở); \textit{nghề nghiệp và thu nhập} (nghề nghiệp, loại hình tổ chức làm
việc, thâm niên công việc, thu nhập khai báo); và \textit{quan hệ với ngân hàng}
(thời gian quan hệ, việc có tài khoản nhận lương, số tài khoản đang mở, tổng số
dư bình quân ba tháng). Danh mục đầy đủ kèm mô tả từng đặc trưng được liệt kê tại
Phụ lục~\ref{app:features}.

Mảng thứ ba phản ánh \textit{mức độ gắn bó} của khách hàng với tổ chức tín dụng,
yếu tố có giá trị dự báo nhưng không xuất hiện trong bộ quy tắc thẩm định truyền
thống.

\subsection{Trích xuất nhóm đặc trưng lịch sử tín dụng và quan hệ CIC}
\label{subsec:dac-trung-cic}

Nhóm này gồm mười sáu đặc trưng lấy từ bản tin CIC, phản ánh bốn khía cạnh của
quan hệ tín dụng: \textit{mức độ tín nhiệm} (điểm CIC), \textit{lịch sử trả nợ}
(nhóm nợ hiện tại, nhóm nợ xấu nhất 12 tháng, số ngày quá hạn tối đa, nợ đã xử lý
rủi ro, số tổ chức từng phát sinh nợ xấu), \textit{mức độ và cơ cấu dư nợ} (số tổ
chức đang có quan hệ, số hợp đồng còn hiệu lực và đã tất toán, tổng dư nợ, tỷ lệ
sử dụng hạn mức, nghĩa vụ trả nợ hàng tháng, độ dài lịch sử tín dụng), và
\textit{mức độ tìm kiếm tín dụng} (số lần truy vấn trong 3 tháng).

Bổ sung vào nhóm là hai biến chỉ báo tạo ra ở bước xử lý khuyết, phản ánh hai
tình huống khác nhau về bản chất: khách hàng \textit{không có} lịch sử tín dụng
và trường hợp \textit{không truy vấn được} dữ liệu.

\subsection{Trích xuất nhóm đặc trưng hành vi và biến tương tác}
\label{subsec:dac-trung-hanh-vi}

\subsubsection{Đặc trưng hành vi theo cửa sổ thời gian}

Từ lịch sử giao dịch, đề án tính hai mươi ba độ đo trên mỗi cửa sổ thời gian
$W \in \{7, 30, 90, 180\}$ ngày, thuộc bốn mảng: \textit{quy mô hoạt động} (số
giao dịch, dòng tiền vào, dòng tiền ra, dòng tiền thuần, giá trị giao dịch trung
bình); \textit{mức biến động} (độ lệch chuẩn và hệ số biến thiên của giá trị giao
dịch); \textit{cơ cấu dòng tiền theo mục đích} (nhận lương, trả nợ, chi tiêu nhóm
rủi ro, rút tiền mặt, thương mại điện tử, hóa đơn tiện ích); và \textit{dấu hiệu
hành vi} (giao dịch ban đêm, giao dịch giá trị lớn, giao dịch qua kênh di động),
cùng năm tỷ số phái sinh từ các độ đo trên.

Bổ sung thêm bảy đặc trưng không phụ thuộc cửa sổ: số ngày kể từ giao dịch gần
nhất, ba biến xu thế so sánh cửa sổ ngắn với cửa sổ dài, tỷ lệ dòng tiền vào thực
tế so với thu nhập khai báo, mức độ đều đặn của việc nhận lương và độ ổn định thu
nhập.

Trong số này, biến \textit{tỷ lệ dòng tiền vào thực tế so với thu nhập khai báo}
có ý nghĩa nghiệp vụ riêng: nó cung cấp một phép kiểm chứng độc lập cho thông tin
thu nhập khách hàng tự khai, công việc mà quy trình thẩm định truyền thống phải
thực hiện thủ công qua đối chiếu chứng từ.

\subsubsection{Thuật toán bảo đảm tính đúng đắn theo thời điểm}
\label{subsubsec:point-in-time}

Yêu cầu tính đúng đắn theo thời điểm đặt ra một thách thức về tính toán. Cách làm
trực tiếp, tức với mỗi hồ sơ lọc toàn bộ giao dịch của khách hàng trong cửa sổ rồi
tổng hợp -- có độ phức tạp $O(N \times |W| \times \bar{m})$ với $N$ là số hồ sơ,
$|W|$ là số cửa sổ và $\bar{m}$ là số giao dịch trung bình mỗi khách hàng. Với
hàng chục nghìn hồ sơ và hàng triệu giao dịch, chi phí này không chấp nhận được.

Đề án sử dụng kỹ thuật \textbf{luỹ kế theo nhóm kết hợp ghép gần nhất theo thời
gian}. Với mỗi độ đo $m$, gọi $C_m(i, t)$ là tổng luỹ kế của độ đo đó cho khách
hàng $i$ tính đến thời điểm $t$. Khi đó tổng của độ đo trong cửa sổ $(t - W, t]$
được tính bằng hiệu hai giá trị luỹ kế:

\begin{equation}
S_m(i, t, W) = C_m(i, t) - C_m(i, t - W).
\label{eq:cumsum-window}
\end{equation}

Hai giá trị luỹ kế được truy xuất bằng phép ghép gần nhất hướng lùi
(\code{merge\_asof}), tương đương thao tác ``tìm bản ghi gần nhất không vượt quá
mốc thời gian'' với độ phức tạp tuyến tính sau khi dữ liệu đã được sắp xếp. Nhờ
đó, toàn bộ đặc trưng cửa sổ cho $N$ hồ sơ được tính bằng $2 \times |W|$ phép
ghép vectơ hóa thay vì $N \times |W|$ truy vấn riêng lẻ.

Phương pháp này còn có một ưu điểm về mặt đúng đắn: vì công thức chỉ sử dụng giá
trị luỹ kế tại hai mốc thời gian, nó \textit{không thể} vô tình bao gồm giao dịch
phát sinh sau thời điểm nộp hồ sơ. Tính đúng đắn theo thời điểm được bảo đảm bởi
chính cấu trúc thuật toán chứ không phụ thuộc vào việc kiểm tra bổ sung.

\subsubsection{Biến tổ hợp và biến tương tác}

Mười ba biến phái sinh được xây dựng nhằm biểu diễn các chỉ tiêu nghiệp vụ then
chốt mà không biến gốc nào thể hiện được:

\begin{table}[htbp]
\centering
\caption{Các biến tổ hợp và biến tương tác}
\label{tab:derived-features}
\small
\setlength{\tabcolsep}{4pt}
\begin{tabular}{@{}L{3.3cm}C{5.3cm}L{6.1cm}@{}}
\toprule
\textbf{Biến} & \textbf{Công thức} & \textbf{Ý nghĩa nghiệp vụ} \\
\midrule
\code{dti} &
$\dfrac{\text{Trả góp mới} + \text{Nợ hiện hữu}}{\text{Thu nhập}}$ &
Tổng gánh nặng nợ trên thu nhập, chỉ tiêu trung tâm của đánh giá khả năng trả nợ
\\[6pt]
\code{pti} &
$\dfrac{\text{trả góp khoản vay mới}}{\text{thu nhập}}$ &
Gánh nặng riêng của khoản vay đang xét \\[6pt]
\code{disposable\_income} &
thu nhập $-$ trả nợ $-$ chi phí sinh hoạt &
Thu nhập còn lại thực tế sau khi trừ nghĩa vụ và chi phí tối thiểu \\[6pt]
\code{credit\_history\_ratio} &
$\dfrac{\text{độ dài lịch sử tín dụng}}{\text{số tháng từ tuổi 18}}$ &
Mức độ tham gia thị trường tín dụng so với độ tuổi \\[6pt]
\code{dpd\_x\_institutions} &
số ngày quá hạn $\times$ số tổ chức &
Biến tương tác: quá hạn tại nhiều tổ chức nghiêm trọng hơn quá hạn tại một tổ
chức \\[6pt]
\code{high\_dti\_short\_tenure} &
$\mathbb{1}[\DTI > 0{,}5 \wedge \text{thâm niên} < 12]$ &
Cờ rủi ro kết hợp: gánh nặng nợ cao đi kèm công việc chưa ổn định \\[6pt]
\code{income\_verification\_ratio} &
$\dfrac{\text{dòng tiền vào 30 ngày}}{\text{thu nhập khai báo}}$ &
Kiểm chứng độc lập thông tin thu nhập tự khai \\
\bottomrule
\end{tabular}
\end{table}

Biến \code{disposable\_income} sử dụng mô hình chi phí sinh hoạt gồm hai thành
phần: một mức sàn cố định phụ thuộc số người phụ thuộc, và một phần tỷ lệ thuận
với thu nhập. Cách mô hình hóa này khắc phục nhược điểm của việc dùng một mức cố
định duy nhất, vốn quá khắt khe với nhóm thu nhập thấp và quá nới lỏng với nhóm
thu nhập cao.

\subsection{Đánh giá mức độ quan trọng của các đặc trưng}
\label{subsec:danh-gia-dac-trung}

Đề án sử dụng ba phương pháp bổ trợ nhau để đánh giá tầm quan trọng của đặc
trưng, mỗi phương pháp trả lời một câu hỏi khác nhau.

\paragraph{Information Value -- năng lực dự báo đơn biến.} Trả lời câu hỏi: mỗi
đặc trưng, xét riêng lẻ, phân biệt hồ sơ tốt và xấu tốt đến mức nào? Ưu điểm là
chi phí tính toán thấp và không phụ thuộc mô hình. Hạn chế là không tính đến
tương tác: một biến có $\IV$ thấp vẫn có thể có giá trị khi kết hợp với biến
khác.

\paragraph{SHAP -- đóng góp trong mô hình đã huấn luyện.} Trả lời câu hỏi: trong
mô hình cụ thể đã được huấn luyện, đặc trưng nào thực sự chi phối kết quả dự báo?
Phương pháp này tính đến tương tác và cho kết quả ở cả cấp độ toàn cục lẫn cấp độ
từng hồ sơ.

\paragraph{Phép hoán vị -- kiểm chứng chéo.} Trả lời câu hỏi: nếu phá vỡ quan hệ
giữa một đặc trưng và nhãn bằng cách xáo trộn ngẫu nhiên giá trị của nó, chất
lượng dự báo suy giảm bao nhiêu? Phương pháp này độc lập với cấu trúc nội bộ của
mô hình, nên dùng làm phép kiểm chứng chéo cho kết quả SHAP.

\subsection{Lựa chọn tập đặc trưng tối ưu cho mô hình}
\label{subsec:lua-chon-dac-trung}

Tập đặc trưng ban đầu gồm \NFeaturesInitial{} biến, trong đó nhiều biến được sinh
theo cùng một mẫu trên bốn cửa sổ thời gian khác nhau nên tương quan mạnh với
nhau. Đề án áp dụng quy trình lọc ba bước nối tiếp, sắp xếp từ rẻ đến đắt về chi
phí tính toán (Hình~\ref{fig:feature-selection}).

\begin{figure}[htbp]
\centering
\begin{tikzpicture}[
  node distance=0.62cm,
  fstep/.style={draw, rounded corners=2pt, align=center, minimum height=1.0cm,
               text width=3.05cm, font=\small, fill=blue!4},
  arrow/.style={-{Stealth[length=2mm]}, thick, draw=black!65},
  note/.style={font=\footnotesize\itshape, align=center, text width=3.05cm},
]
\node[fstep] (s0) {Tập đặc trưng\\ban đầu};
\node[fstep, right=of s0] (s1) {Bước 1\\Lọc theo $\IV$};
\node[fstep, right=of s1] (s2) {Bước 2\\Lọc theo tương quan};
\node[fstep, right=of s2] (s3) {Bước 3\\Lọc theo $\PSI$};

\draw[arrow] (s0) -- (s1);
\draw[arrow] (s1) -- (s2);
\draw[arrow] (s2) -- (s3);

\node[note, below=0.32cm of s1] {loại biến\\không phân biệt được};
\node[note, below=0.32cm of s2] {loại biến\\dư thừa};
\node[note, below=0.32cm of s3] {loại biến\\không ổn định};
\end{tikzpicture}
\caption{Quy trình lựa chọn tập đặc trưng qua ba bộ lọc nối tiếp}
\label{fig:feature-selection}
\end{figure}

\paragraph{Bước 1 -- lọc theo năng lực dự báo.} Loại các đặc trưng có
$\IV < 0{,}02$, ngưỡng mà theo thực hành ngành tương ứng với ``không có khả năng
phân biệt''. Ngưỡng trên được đặt ở mức $\IV = 2{,}0$ để phát hiện trường hợp rò
rỉ nhãn. Lưu ý rằng nhóm đặc trưng CIC vốn có $\IV$ cao (khoảng 0,9--1,3) do phản
ánh trực tiếp lịch sử trả nợ; các biến này đã được kiểm tra tính đúng đắn theo
thời điểm nên $\IV$ cao là bản chất của dữ liệu chứ không phải dấu hiệu rò rỉ.

\paragraph{Bước 2 -- lọc theo tương quan.} Trong mỗi cặp biến có hệ số tương quan
tuyệt đối vượt 0,95, giữ lại biến có $\IV$ cao hơn. Bước này chủ yếu loại bỏ các
biến trùng lặp giữa các cửa sổ thời gian liền kề.

\paragraph{Bước 3 -- lọc theo độ ổn định.} Loại các đặc trưng có $\PSI > 0{,}25$
giữa giai đoạn huấn luyện và giai đoạn kiểm thử. Mô hình phụ thuộc vào biến không
ổn định sẽ suy giảm nhanh khi vận hành.

\section{Xây dựng nhãn và xác định quy tắc phê duyệt}
\label{sec:xay-dung-nhan}

\subsection{Định nghĩa trạng thái rủi ro và tiêu chuẩn nợ xấu}
\label{subsec:dinh-nghia-nhan}

\paragraph{Tiêu chuẩn nợ xấu.} Đề án sử dụng định nghĩa theo quy định phân loại
nợ hiện hành~\cite{tt11}: khoản vay bị coi là nợ xấu khi được phân vào nhóm 3, 4
hoặc 5, tương ứng với số ngày quá hạn từ 90 ngày trở lên. Định nghĩa này trùng
với ngưỡng mà chuẩn mực Basel~\cite{bcbs2006} sử dụng cho khái niệm vỡ nợ, nhờ đó
kết quả nghiên cứu có thể đối chiếu với các công trình quốc tế.

\paragraph{Cửa sổ quan sát.} Một hồ sơ chỉ được gán nhãn sau khi đã trải qua cửa
sổ quan sát hiệu quả trả nợ đủ dài. Đề án chọn cửa sổ 12 tháng dựa trên hai căn
cứ đối lập nhau cần cân bằng: cửa sổ quá ngắn bỏ sót phần lớn trường hợp vỡ nợ
muộn, làm tỷ lệ nợ xấu quan sát được thấp hơn thực tế và mô hình lệch về phía lạc
quan; cửa sổ quá dài buộc phải loại nhiều hồ sơ gần đây khỏi tập huấn luyện, làm
dữ liệu cũ đi và mô hình kém phản ánh điều kiện hiện tại.

Căn cứ định lượng cho lựa chọn này là \textbf{phân tích vintage}: đo tỷ lệ nợ xấu
tích luỹ theo từng tháng kể từ thời điểm cấp tín dụng. Điểm mà đường cong tích
luỹ bắt đầu đi ngang chính là độ dài cửa sổ hợp lý. Kết quả phân tích được trình
bày ở mục~\ref{subsec:kq-gan-nhan}.

\subsection{Phương pháp gán nhãn dữ liệu hồ sơ vay}
\label{subsec:phuong-phap-gan-nhan}

Quy trình gán nhãn gồm ba bước:

\begin{enumerate}[label=(\arabic*)]
  \item \textbf{Xác định thời điểm kết thúc cửa sổ quan sát} cho từng hồ sơ, bằng
  thời điểm cấp tín dụng cộng với độ dài cửa sổ.

  \item \textbf{Phân loại trạng thái.} Hồ sơ có thời điểm kết thúc cửa sổ vượt
  quá mốc trích xuất dữ liệu được gán \code{INDETERMINATE}. Trong số các hồ sơ
  còn lại, hồ sơ phát sinh nợ xấu \textit{trong} cửa sổ được gán \code{BAD}, phần
  còn lại được gán \code{GOOD}. Trường hợp vỡ nợ muộn hơn cửa sổ không được tính
  là \code{BAD}, vì nằm ngoài phạm vi dự báo của mô hình.

  \item \textbf{Loại nhóm chưa xác định} khỏi tập dữ liệu mô hình hóa.
\end{enumerate}

\subsection{Xác định điểm cắt phê duyệt}
\label{subsec:diem-cat}

\subsubsection{Vì sao ngưỡng 0,5 không phù hợp}

Ngưỡng mặc định 0,5 của bài toán phân loại nhị phân không phù hợp với bài toán
tín dụng vì hai lý do. Thứ nhất, dữ liệu mất cân bằng: với tỷ lệ nợ xấu khoảng
15\%, rất ít hồ sơ có $\PD$ vượt 0,5, nên dùng ngưỡng này gần như duyệt toàn bộ.
Thứ hai, chi phí hai loại sai lầm rất khác nhau: bỏ sót một hồ sơ xấu gây tổn
thất bằng một phần đáng kể dư nợ, trong khi từ chối sai một hồ sơ tốt chỉ làm mất
phần lãi biên.

\subsubsection{Ba tiêu chí xác định điểm cắt}

Đề án hiện thực ba phương pháp xác định điểm cắt, tương ứng ba cách đặt bài toán
khác nhau trong thực tế nghiệp vụ.

\paragraph{Tiêu chí 1 -- tối thiểu chi phí kỳ vọng.} Chọn ngưỡng $\tau$ làm cực
tiểu chi phí kỳ vọng:

\begin{equation}
\tau^* = \argmin_{\tau}
\left[ c_{FN} \cdot \text{FN}(\tau) + c_{FP} \cdot \text{FP}(\tau) \right] / n,
\label{eq:cost-threshold}
\end{equation}

với $c_{FN} : c_{FP} = 12 : 1$. Áp dụng khi tổ chức tín dụng ước lượng được chi
phí của hai loại sai lầm.

\paragraph{Tiêu chí 2 -- đạt tỷ lệ duyệt mục tiêu.} Chọn ngưỡng bằng phân vị
tương ứng với tỷ lệ duyệt mong muốn. Áp dụng khi ngưỡng bị ràng buộc bởi kế hoạch
tăng trưởng dư nợ.

\paragraph{Tiêu chí 3 -- bảo đảm trần nợ xấu.} Chọn ngưỡng cao nhất mà tỷ lệ nợ
xấu trong nhóm được duyệt vẫn dưới mức trần cho phép. Chọn ngưỡng \textit{cao
nhất} thỏa mãn ràng buộc nhằm tối đa hóa quy mô cấp tín dụng trong giới hạn khẩu
vị rủi ro.

\subsubsection{Hai ngưỡng cho ba vùng quyết định}

Hệ thống hỗ trợ kiểm duyệt cần \textbf{hai} ngưỡng chứ không phải một, chia thang
$\PD$ thành ba vùng:

\begin{equation}
\text{Khuyến nghị} =
\begin{cases}
\texttt{APPROVE}, & \PD \le \tau_{\text{duyệt}}, \\
\texttt{REVIEW},  & \tau_{\text{duyệt}} < \PD \le \tau_{\text{từ chối}}, \\
\texttt{REJECT},  & \PD > \tau_{\text{từ chối}}.
\end{cases}
\label{eq:three-zones}
\end{equation}

Hai ngưỡng được đặt đối xứng quanh điểm cắt tối ưu chi phí, sao cho tỷ lệ hồ sơ
rơi vào vùng chuyển thẩm định xấp xỉ một mức mục tiêu định trước. Tỷ lệ này là
\textit{tham số chính sách} chứ không phải tham số kỹ thuật: tăng lên thì mức độ
an toàn cao hơn nhưng khối lượng thẩm định thủ công lớn hơn; giảm xuống thì tự
động hóa cao hơn nhưng rủi ro khuyến nghị sai tăng. Đề án đặt mục tiêu vùng thẩm
định ở mức 20\%.

\section{Thiết kế mô hình học máy kiểm duyệt tín dụng}
\label{sec:thiet-ke-mo-hinh}

\subsection{Lựa chọn các thuật toán học máy dự kiến}
\label{subsec:lua-chon-thuat-toan}

Đề án lựa chọn mười hai mô hình để so sánh, trong đó năm thuật toán là trọng tâm
nghiên cứu -- hồi quy logistic, cây quyết định, rừng ngẫu nhiên, máy vectơ hỗ trợ
và gradient boosting -- và bảy mô hình còn lại là các biến thể hiện đại của cùng
những họ thuật toán đó, đưa vào nhằm đánh giá xem mức cải thiện thực tế có tương
xứng với chi phí triển khai hay không
(Bảng~\ref{tab:model-selection}).

\begin{table}[htbp]
\centering
\caption{Danh mục mô hình được thiết kế và so sánh}
\label{tab:model-selection}
\small
\begin{tabular}{@{}clL{3.0cm}L{5.2cm}@{}}
\toprule
\textbf{TT} & \textbf{Mô hình} & \textbf{Nhóm} & \textbf{Vai trò trong so sánh} \\
\midrule
1  & Hồi quy logistic      & Tuyến tính & Mô hình cơ sở tham chiếu \\
2  & Thẻ điểm WOE          & Tuyến tính & Chuẩn mực thực tiễn ngân hàng \\
3  & Cây quyết định (CART) & Cây đơn    & Mô hình phi tuyến đơn giản nhất \\
4  & Rừng ngẫu nhiên       & Bagging    & Học tập hợp song song \\
5  & Extra Trees           & Bagging    & Biến thể tăng mức ngẫu nhiên hóa \\
6  & Máy vectơ hỗ trợ      & SVM        & Tiếp cận theo lề cực đại \\
7  & Hist Gradient Boosting& Boosting   & Học tập hợp tuần tự \\
8  & XGBoost               & Boosting   & Biến thể có điều chuẩn tường minh \\
9  & LightGBM              & Boosting   & Biến thể sinh cây theo chiều lá \\
10 & CatBoost              & Boosting   & Biến thể xử lý biến định tính \\
11 & Mạng nơ-ron (MLP)     & Mạng nơ-ron& Mô hình phi tuyến toàn cục \\
12 & Kết hợp xếp tầng      & Ensemble   & Tổ hợp các mô hình cơ sở khác biệt \\
\bottomrule
\end{tabular}
\end{table}

\paragraph{Ghi chú về hiện thực SVM.} Như đã phân tích ở mục~\ref{subsec:svm},
SVM hạt nhân có chi phí huấn luyện từ $O(n^2)$ đến $O(n^3)$, không khả thi với
tập huấn luyện hàng chục nghìn hồ sơ. Đề án sử dụng phép xấp xỉ Nystroem với 600
điểm mốc, sau đó huấn luyện SVM tuyến tính trên không gian đặc trưng xấp xỉ. Xác
suất đầu ra được quy đổi từ giá trị hàm quyết định bằng phương pháp Platt qua
kiểm định chéo ba tầng.

\subsection{Thiết kế kiến trúc tổng thể của mô hình học máy hỗ trợ ra quyết định}
\label{subsec:kien-truc-tong-the}

\subsubsection{Kiến trúc phân tầng}

Hệ thống được tổ chức thành bốn tầng (Hình~\ref{fig:architecture}), tách bạch
giữa phần \textit{xây dựng mô hình} (thực hiện định kỳ, ngoại tuyến) và phần
\textit{sử dụng mô hình} (thực hiện mỗi khi có hồ sơ mới).

\begin{figure}[htbp]
\centering
\begin{tikzpicture}[
  node distance=0.5cm and 0.5cm,
  layer/.style={draw, rounded corners=3pt, align=center, font=\small\bfseries,
                minimum height=0.8cm, text width=2.0cm},
  comp/.style={draw, rounded corners=2pt, align=center, font=\footnotesize,
               minimum height=0.85cm, text width=2.4cm, fill=blue!4},
  arrow/.style={-{Stealth[length=2mm]}, thick, draw=black!65},
]
% Tầng 1: dữ liệu
\node[layer, fill=gray!12] (l1) {Tầng\\dữ liệu};
\node[comp, right=of l1] (d1) {Hồ sơ khách hàng\\và tài khoản};
\node[comp, right=of d1] (d2) {Đơn đề nghị vay};
\node[comp, right=of d2] (d3) {Bản tin CIC};
\node[comp, right=of d3] (d4) {Lịch sử giao dịch};

% Tầng 2: xử lý
\node[layer, fill=gray!12, below=1.0cm of l1] (l2) {Tầng\\xử lý dữ liệu};
\node[comp, right=of l2, fill=orange!8] (p1) {Kiểm tra\\chất lượng};
\node[comp, right=of p1, fill=orange!8] (p2) {Tiền xử lý\\và làm sạch};
\node[comp, right=of p2, fill=orange!8] (p3) {Trích xuất\\đặc trưng};
\node[comp, right=of p3, fill=orange!8] (p4) {Lựa chọn\\đặc trưng};

% Tầng 3: mô hình
\node[layer, fill=gray!12, below=1.0cm of l2] (l3) {Tầng\\mô hình};
\node[comp, right=of l3, fill=green!8] (m1) {Huấn luyện\\và tối ưu};
\node[comp, right=of m1, fill=green!8] (m2) {Hiệu chuẩn\\xác suất};
\node[comp, right=of m2, fill=green!8] (m3) {Giải thích\\(SHAP)};
\node[comp, right=of m3, fill=green!8] (m4) {Giám sát\\trôi dạt};

% Tầng 4: quyết định
\node[layer, fill=gray!12, below=1.0cm of l3] (l4) {Tầng hỗ trợ\\quyết định};
\node[comp, right=of l4, fill=red!6] (e1) {Luật\\nghiệp vụ};
\node[comp, right=of e1, fill=red!6] (e2) {Chấm điểm\\$\PD$};
\node[comp, right=of e2, fill=red!6] (e3) {Chính sách\\hạn mức};
\node[comp, right=of e3, fill=red!6] (e4) {Khuyến nghị\\và lý do};

\draw[arrow] (d2.south) -- ++(0,-0.32) -| (p1.north);
\draw[arrow] (p1) -- (p2);
\draw[arrow] (p2) -- (p3);
\draw[arrow] (p3) -- (p4);
\draw[arrow] (p3.south) -- ++(0,-0.32) -| (m1.north);
\draw[arrow] (m1) -- (m2);
\draw[arrow] (m2.south) -- ++(0,-0.32) -| (e2.north);
\draw[arrow] (e1) -- (e2);
\draw[arrow] (e2) -- (e3);
\draw[arrow] (e3) -- (e4);
\end{tikzpicture}
\caption{Kiến trúc phân tầng của mô hình hỗ trợ kiểm duyệt tín dụng}
\label{fig:architecture}
\end{figure}

\subsubsection{Logic tổng hợp khuyến nghị}

Khuyến nghị cuối cùng là kết quả của một chuỗi ưu tiên rõ ràng, trong đó luật
nghiệp vụ có quyền phủ quyết cao hơn mô hình thống kê:

\begin{enumerate}[label=\arabic*.]
  \item Có luật loại trừ cứng bị vi phạm $\Rightarrow$ \code{REJECT}
  \item Mô hình gặp lỗi khi chấm điểm $\Rightarrow$ \code{REVIEW}
  \item Không truy vấn được dữ liệu CIC $\Rightarrow$ \code{REVIEW}
  \item $\PD > \tau_{\text{từ chối}}$ $\Rightarrow$ \code{REJECT}
  \item Có luật yêu cầu thẩm định bị vi phạm $\Rightarrow$ \code{REVIEW}
  \item $\PD > \tau_{\text{duyệt}}$ $\Rightarrow$ \code{REVIEW}
  \item Hạn mức tính được bằng 0 $\Rightarrow$ \code{REJECT}
  \item Các trường hợp còn lại $\Rightarrow$ \code{APPROVE}
\end{enumerate}

Thứ tự này phản ánh một nguyên tắc thiết kế: \textit{ràng buộc tuân thủ đứng
trước đánh giá thống kê}. Một hồ sơ vi phạm điều kiện pháp lý hoặc chính sách
phải bị từ chối bất kể mô hình cho điểm rủi ro thấp đến đâu. Ngược lại, khi hệ
thống không đủ thông tin hoặc không đủ chắc chắn, kết quả mặc định là chuyển cho
con người xem xét chứ không phải tự quyết.

\subsubsection{Bộ luật nghiệp vụ}

Đề án xây dựng bộ luật gồm hai mươi mốt luật, chia theo sáu nhóm và ba mức hành
động (Bảng~\ref{tab:rule-groups}).

\begin{table}[htbp]
\centering
\caption{Cơ cấu bộ luật nghiệp vụ theo nhóm và mức hành động}
\label{tab:rule-groups}
\small
\begin{tabular}{@{}L{3.2cm}L{4.2cm}L{5.0cm}@{}}
\toprule
\textbf{Nhóm luật} & \textbf{Nội dung kiểm tra} & \textbf{Ví dụ} \\
\midrule
Điều kiện vay & Độ tuổi, thu nhập tối thiểu, thâm niên công việc &
Tuổi ngoài khoảng 20--62 $\Rightarrow$ từ chối \\
Lịch sử tín dụng & Nhóm nợ, số ngày quá hạn, nợ đã xử lý rủi ro &
Nhóm nợ xấu nhất 12 tháng $> 2$ $\Rightarrow$ từ chối \\
Khả năng trả nợ & $\DTI$, thu nhập còn lại sau trả nợ &
$\DTI > 0{,}6$ $\Rightarrow$ từ chối \\
Tài sản bảo đảm & Tỷ lệ vay trên giá trị tài sản &
$\LTV > 0{,}85$ $\Rightarrow$ chuyển thẩm định \\
Chất lượng dữ liệu & Thiếu dữ liệu CIC, khách hàng thin-file &
Không có bản tin CIC $\Rightarrow$ chuyển thẩm định \\
Dấu hiệu bất thường & Số lần truy vấn tín dụng, hành vi giao dịch &
Truy vấn tín dụng $> 5$ lần trong 3 tháng $\Rightarrow$ chuyển thẩm định \\
\bottomrule
\end{tabular}
\end{table}

Ba mức hành động gồm: \code{HARD\_REJECT} (từ chối, có quyền phủ quyết cao nhất),
\code{REVIEW} (buộc chuyển thẩm định thủ công) và \code{WARNING} (ghi nhận vào
danh sách lý do nhưng không thay đổi khuyến nghị).

\subsection{Xác định các hàm mất mát}
\label{subsec:ham-mat-mat}

\paragraph{Hàm mất mát huấn luyện.} Toàn bộ mô hình sử dụng hàm entropy chéo nhị
phân~\eqref{eq:logloss}. Lựa chọn này không phải quy ước tùy tiện: log-loss là
quy tắc chấm điểm chính đáng, nghĩa là kỳ vọng của nó đạt cực tiểu khi và chỉ khi
giá trị dự báo bằng đúng xác suất thật. Đặc điểm này có ý nghĩa quyết định vì hệ
thống sử dụng giá trị $\PD$ trực tiếp chứ không chỉ dùng để xếp hạng.

\paragraph{Bù mất cân bằng lớp.} Với tỷ lệ nợ xấu khoảng 15\%, hàm mất mát được
bổ sung trọng số lớp:

\begin{equation}
\ell_w(y, p) = -\big[\, w_1 \, y \log p + w_0 \,(1 - y) \log (1 - p) \,\big],
\qquad
w_1 = \frac{n_0}{n_1},\ w_0 = 1,
\label{eq:weighted-logloss}
\end{equation}

trong đó $n_0, n_1$ là số quan sát của lớp âm và lớp dương. Việc dùng trọng số
lớp làm xác suất đầu ra bị lệch hệ thống về phía cao hơn thực tế; đây là lý do
bước hiệu chuẩn xác suất trở thành bắt buộc chứ không phải tùy chọn.

\paragraph{Hàm chi phí nghiệp vụ.} Tách biệt với hàm mất mát huấn luyện, hàm chi
phí nghiệp vụ được dùng ở bước chọn ngưỡng quyết định:

\begin{equation}
\mathcal{C}(\tau) = \frac{1}{n}\Big[ c_{FN} \cdot \text{FN}(\tau)
                                   + c_{FP} \cdot \text{FP}(\tau) \Big].
\label{eq:business-cost}
\end{equation}

Việc tách hai hàm này là một lựa chọn thiết kế có chủ đích: mô hình được huấn
luyện để ước lượng xác suất đúng, còn chính sách chi phí được áp dụng ở bước
quyết định. Nhờ vậy, khi tổ chức tín dụng thay đổi khẩu vị rủi ro, chỉ cần điều
chỉnh ngưỡng mà không phải huấn luyện lại mô hình.

\subsection{Lựa chọn chỉ số đánh giá hiệu năng}
\label{subsec:chi-so-danh-gia}

Trên cơ sở bộ chỉ tiêu đã trình bày ở mục~\ref{subsec:so-sanh-mo-hinh}, đề
án bổ sung nhóm chỉ tiêu thứ năm mang tính nghiệp vụ nhằm quy đổi chất lượng
thống kê thành hệ quả kinh doanh. Bảng~\ref{tab:metrics-design} nêu vai trò của
từng nhóm trong quy trình đánh giá.

\begin{table}[htbp]
\centering
\caption{Hệ thống chỉ tiêu đánh giá và vai trò trong quy trình}
\label{tab:metrics-design}
\small
\begin{tabular}{@{}L{2.4cm}L{3.4cm}L{6.6cm}@{}}
\toprule
\textbf{Nhóm} & \textbf{Chỉ tiêu} & \textbf{Vai trò} \\
\midrule
Phân loại phổ quát & Accuracy, Precision, Recall, F1, PR-AUC &
Báo cáo đầy đủ; Recall được ưu tiên hơn Precision do chi phí bỏ sót hồ sơ xấu cao
\\[2pt]
Đặc thù tín dụng & $\AUC$, Gini, $\KS$ &
Tiêu chí chính để xếp hạng và lựa chọn mô hình; không phụ thuộc ngưỡng \\[2pt]
Hiệu chuẩn & Brier, ECE, đường cong hiệu chuẩn &
Điều kiện cần để dùng $\PD$ trong phân vùng quyết định và tính tổn thất kỳ vọng
\\[2pt]
Vận hành & Thời gian huấn luyện, thời gian chấm điểm, mức độ giải thích được &
Quyết định khả năng triển khai thực tế \\[2pt]
Nghiệp vụ & Tỷ lệ duyệt, tỷ lệ nợ xấu nhóm duyệt, chi phí kỳ vọng, bảng gains &
Quy đổi chất lượng thống kê thành hệ quả kinh doanh \\
\bottomrule
\end{tabular}
\end{table}

\subsection{Tích hợp cơ chế giải thích mô hình}
\label{subsec:tich-hop-giai-thich}

Cơ chế giải thích được tích hợp ở hai cấp độ.

\paragraph{Cấp độ toàn cục.} Phân tích SHAP trên một mẫu đại diện của tập kiểm
thử cho biết đặc trưng nào chi phối kết quả dự báo của mô hình nói chung, và mỗi
nhóm nguồn dữ liệu đóng góp bao nhiêu phần vào quyết định. Kết quả này phục vụ
hai mục đích: kiểm chứng rằng mô hình học được quan hệ hợp lý về mặt nghiệp vụ,
và định lượng giá trị của từng nguồn dữ liệu để làm cơ sở cho quyết định đầu tư
dữ liệu.

\paragraph{Cấp độ từng hồ sơ.} Với mỗi hồ sơ được chấm điểm, mô hình sinh danh
sách lý do gồm ba thành phần theo thứ tự ưu tiên: (i) các luật nghiệp vụ bị vi
phạm -- đây là loại lý do tường minh và dễ hiểu nhất; (ii) các chỉ tiêu định
lượng then chốt kèm giá trị thực tế của hồ sơ; và (iii) các đặc trưng có đóng góp
SHAP lớn nhất. Danh sách này được trả về cùng kết quả chấm điểm, cho phép cán bộ
tín dụng đánh giá xem khuyến nghị của mô hình có hợp lý hay không.

\section{Quy trình huấn luyện và tối ưu mô hình}
\label{sec:quy-trinh-huan-luyen}

\subsection{Thiết lập môi trường và công cụ huấn luyện}
\label{subsec:moi-truong-huan-luyen}

Toàn bộ quy trình được hiện thực bằng Python, tổ chức thành một gói phần mềm mà
mỗi gói con tương ứng một mục của chương này: \code{data} cho khâu thu thập và
kiểm tra chất lượng dữ liệu, \code{preprocessing} cho tiền xử lý, \code{features}
cho trích xuất và lựa chọn đặc trưng, \code{labeling} cho gán nhãn và điểm cắt,
\code{models} cho huấn luyện và đánh giá mô hình, \code{decision} cho bộ máy hỗ
trợ quyết định, \code{serving} cho dịch vụ REST và \code{monitoring} cho giám
sát. Cấu trúc thư mục chi tiết được trình bày tại Phụ lục~\ref{app:source}.

Các thư viện chính gồm: \comp{scikit-learn} cho các mô hình cơ sở và khung
pipeline; \comp{XGBoost}, \comp{LightGBM}, \comp{CatBoost} cho các biến thể
gradient boosting; \comp{SHAP} cho giải thích mô hình; \comp{pandas} và
\comp{NumPy} cho xử lý dữ liệu; \comp{FastAPI} cho dịch vụ REST. Toàn bộ tham số
được đặt tập trung trong một file cấu hình duy nhất, và hạt giống ngẫu nhiên được
cố định để bảo đảm khả năng tái lập.

\subsection{Quy trình thử nghiệm và kiểm tra chéo}
\label{subsec:kiem-tra-cheo}

\paragraph{Kiểm định ngoài thời gian cho đánh giá cuối cùng.} Việc đánh giá chất
lượng mô hình được thực hiện trên tập kiểm thử gồm các hồ sơ thuộc giai đoạn sau
giai đoạn huấn luyện. Đây là phương pháp phản ánh đúng điều kiện vận hành.

\paragraph{Kiểm định chéo theo chuỗi thời gian cho tinh chỉnh tham số.} Với việc
tinh chỉnh siêu tham số, đề án sử dụng \code{TimeSeriesSplit} -- phương pháp chia
tập huấn luyện thành nhiều lần khớp liên tiếp, trong đó mỗi lần dùng dữ liệu quá
khứ để huấn luyện và dữ liệu kế tiếp để đánh giá:

\begin{center}
\begin{tabular}{@{}ll@{}}
\toprule
\textbf{Lần khớp} & \textbf{Phân chia} \\
\midrule
1 & Huấn luyện: đoạn 1 \quad$\rightarrow$\quad Đánh giá: đoạn 2 \\
2 & Huấn luyện: đoạn 1--2 \quad$\rightarrow$\quad Đánh giá: đoạn 3 \\
3 & Huấn luyện: đoạn 1--3 \quad$\rightarrow$\quad Đánh giá: đoạn 4 \\
\bottomrule
\end{tabular}
\end{center}

Phương pháp này khác với kiểm định chéo $k$-fold thông thường ở chỗ nó không bao
giờ dùng dữ liệu tương lai để huấn luyện -- điều kiện bắt buộc với dữ liệu có cấu
trúc thời gian.

\subsection{Phương pháp tối ưu hóa siêu tham số}
\label{subsec:toi-uu-sieu-tham-so}

Đề án sử dụng \textbf{tìm kiếm ngẫu nhiên} thay vì tìm kiếm lưới. Lựa chọn này
dựa trên kết quả của Bergstra và Bengio~\cite{bergstra2012}: với cùng một ngân
sách tính toán, tìm kiếm ngẫu nhiên thường tìm được cấu hình tốt hơn tìm kiếm
lưới. Nguyên nhân là trong không gian siêu tham số nhiều chiều, chỉ một số ít
chiều thực sự ảnh hưởng tới kết quả; tìm kiếm lưới lãng phí phần lớn số lần thử
để dò các giá trị khác nhau của những chiều không quan trọng, trong khi tìm kiếm
ngẫu nhiên bao phủ được nhiều giá trị khác biệt hơn ở những chiều quan trọng.

Không gian tìm kiếm được xác định cho ba mô hình có nhiều siêu tham số nhất --
LightGBM, XGBoost và rừng ngẫu nhiên -- bao gồm các tham số về độ sâu/số lá, tốc
độ học, số quan sát tối thiểu tại lá, tỷ lệ lấy mẫu quan sát và đặc trưng, cùng
các hệ số điều chuẩn. Tiêu chí lựa chọn là $\AUC$ trung bình trên các lần khớp
của kiểm định chéo theo chuỗi thời gian.

Để bảo đảm tính nhất quán, việc tinh chỉnh siêu tham số phải được thực hiện trên
\textit{đúng tập đặc trưng} mà mô hình cuối cùng sẽ sử dụng. Nếu
tinh chỉnh trên tập đầy đủ rồi áp dụng tham số cho tập đã rút gọn, siêu tham số
tìm được không còn tối ưu.

\subsection{Mức độ kiểm soát hiện tượng quá khớp}
\label{subsec:kiem-soat-qua-khop}

Đề án áp dụng bốn cơ chế kiểm soát quá khớp, tác động ở các giai đoạn khác nhau
của quy trình.

\paragraph{Điều chuẩn trong hàm mục tiêu.} Hồi quy logistic và SVM sử dụng điều
chuẩn $L_2$; các mô hình boosting sử dụng đồng thời điều chuẩn $L_1$ và $L_2$
trên giá trị lá.

\paragraph{Ràng buộc cấu trúc mô hình.} Giới hạn độ sâu cây, số lá tối đa và số
quan sát tối thiểu tại mỗi lá. Với cây quyết định đơn, bổ sung tỉa nhánh theo độ
phức tạp chi phí~\eqref{eq:ccp}.

\paragraph{Dừng sớm.} Các mô hình boosting theo dõi $\AUC$ trên tập
\code{train\_es} và dừng khi chỉ tiêu này ngừng cải thiện trong một số vòng lặp
liên tiếp. Cần chỉ định $\AUC$ là chỉ tiêu dừng \textit{duy nhất}: khi sử dụng
trọng số lớp, hàm log-loss trên tập đánh giá xấu đi rất sớm do xác suất bị lệch
hệ thống, nên nếu để cơ chế mặc định theo dõi nhiều chỉ tiêu, quá trình huấn
luyện sẽ dừng quá sớm và mô hình bị dưới khớp.

\paragraph{Lấy mẫu ngẫu nhiên.} Rừng ngẫu nhiên và các mô hình boosting sử dụng
lấy mẫu con quan sát và đặc trưng ở mỗi bước, vừa giảm tương quan giữa các cây
vừa đóng vai trò điều chuẩn ngẫu nhiên.

Mức độ quá khớp được đo bằng chênh lệch $\AUC$ giữa tập kiểm định và tập kiểm
thử; chênh lệch lớn là dấu hiệu mô hình phụ thuộc quá mức vào đặc thù của giai
đoạn huấn luyện. Kết quả đo được trình bày tại mục~\ref{subsec:kq-on-dinh}.

\subsection{Phương pháp đánh giá độ ổn định mô hình theo thời gian}
\label{subsec:do-on-dinh}

Độ ổn định được đánh giá trên hai phương diện.

\paragraph{Ổn định của dữ liệu đầu vào.} Đo bằng chỉ số $\PSI$~\eqref{eq:psi} cho
từng đặc trưng, so sánh phân phối giữa giai đoạn huấn luyện và giai đoạn kiểm
thử. Ngưỡng cảnh báo 0,10 và ngưỡng báo động 0,25.

\paragraph{Ổn định của chất lượng dự báo.} Chia tập kiểm thử theo tháng và tính
$\AUC$, $\KS$ cùng tỷ lệ nợ xấu thực tế cho từng tháng. Xu hướng giảm liên tục
của $\AUC$ là dấu hiệu mô hình đang mất dần hiệu lực và cần được huấn luyện lại.

Hai phương diện này bổ trợ nhau: $\PSI$ phát hiện sớm sự thay đổi của dữ liệu đầu
vào ngay cả khi chưa có nhãn, trong khi chất lượng dự báo theo tháng chỉ đo được
sau khi nhãn đã sẵn sàng. Trong vận hành thực tế, $\PSI$ đóng vai trò cảnh báo
sớm còn chất lượng dự báo đóng vai trò xác nhận.

\section*{Kết luận chương 2}
\addcontentsline{toc}{section}{Kết luận chương 2}

Chương 2 đã trình bày quá trình xây dựng mô hình học máy hỗ trợ kiểm duyệt tín
dụng, từ
phân tích bài toán tới thiết kế quy trình huấn luyện mô hình, với bốn nhóm kết
quả.

Về \textit{phân tích bài toán}: xác định biến mục tiêu là xác suất phát sinh nợ
xấu trong 12 tháng, phân biệt ba trạng thái khi gán nhãn, thiết kế thang bảy hạng
rủi ro và năm thành phần của kết quả đầu ra; hệ thống hóa mười yêu cầu chức năng,
bốn yêu cầu phi chức năng và năm nhóm yêu cầu pháp lý.

Về \textit{xử lý dữ liệu}: thiết lập quy tắc ghép nối bảo đảm tính đúng đắn theo
thời điểm, xây dựng bộ mười lăm ràng buộc kiểm tra tính hợp lệ kèm cơ chế kiểm
chứng bằng phép tiêm lỗi, phân biệt hai cách xử lý khác nhau cho dữ liệu không
hợp lệ và giá trị ngoại lệ, thiết kế chiến lược phân chia bốn tập với vai trò
tách bạch.

Về \textit{đặc trưng}: trình bày cách trích xuất năm nhóm đặc trưng cùng thuật
toán luỹ kế bảo đảm tính đúng đắn theo thời điểm với chi phí tuyến tính, và quy
trình lựa chọn ba bước theo năng lực dự báo, mức tương quan và độ ổn định.

Về \textit{mô hình}: lựa chọn mười hai mô hình để so sánh, thiết kế kiến trúc
phân tầng với logic tổng hợp khuyến nghị theo tám mức ưu tiên, tách hàm mất mát
huấn luyện khỏi hàm chi phí nghiệp vụ, và thiết lập bốn cơ chế kiểm soát quá
khớp.

Toàn bộ thiết kế này được kiểm chứng bằng thực nghiệm ở
Chương~\ref{chap:thuc-nghiem}.
